\section*{Chapter \ref{ch:hf:futures-market-making} --- Futures Market Making}

\begin{solution}{exo:hf:futures-market-making:1}
Time priority: 100 and 150 to the first two. Pro rata: $250\times$ each share of 1,000, rounded down: 25, 100, 75, 50. Top order then pro rata: 100
to the top order, then 150 shared among 900 lots: 66, 50, 33 rounded down, and the leftover lot to the oldest, 67.
\end{solution}

\begin{solution}{exo:hf:futures-market-making:2}
The book holds 2,300 lots: an aggressor of 1,000 gives each $1{,}000\times230/2{,}300=100$; one of 7,907 exhausts the book and gives each its full
230.
\end{solution}

\begin{solution}{exo:hf:futures-market-making:3}
Implied spread bid $9{,}612-9{,}599=13$, offer $9{,}613-9{,}598=15$. Buy the spread at 12, sell the front at 9,612 and buy the back at 9,599:
one tick, less three fees.
\end{solution}

\begin{solution}{exo:hf:futures-market-making:4}
Each maker's best response to the others' sizes is to show more, whatever they do, but when all do, the shares are unchanged and every maker bears
more risk from large orders: utility 19.5 at the equilibrium against 22.3 if all showed 100, with ten makers.
\end{solution}

\begin{solution}{exo:hf:futures-market-making:5}
Makers show far more than they want to trade and cancel most of it whenever the risk of a large fill rises; the shown size is a claim on a share,
not an intention to trade it all.
\end{solution}

\begin{solution}{exo:hf:futures-market-making:6}
It rewards improving the price rather than showing size: the first order at a new price is filled first, so makers compete to set the price with
modest size, which moderates the size race at existing prices.
\end{solution}

\begin{solution}{exo:hf:futures-market-making:7}
Under time priority a maker would show 50 lots; ten makers show 73, a factor of 1.5 instead of 2.3: a larger penalty for inventory restrains the
spiral.
\end{solution}

\begin{solution}{exo:hf:futures-market-making:8}
Most of the 20,000 lots is \omterm{def:hf:futures-market-making:over}{over-quoting} that will be cancelled as soon as a large order looks likely; the depth that stays for a 500-lot order is
what makers want to trade, a fraction of what is shown.
\end{solution}

\begin{solution}{pb:hf:futures-market-making:1}
\begin{enumerate}
\item See Book 1, chapter 19 and the dated box: pro-rated shares rounded down, excess lots first in first out, a top order first where included.
\item 100/150; 25/100/75/50; 100/67/50/33.
\item See \cref{def:hf:futures-market-making:over}.
\item Spreads at one tick nearly always, depth two orders of magnitude above the mean market order, cancellation rates above 96\%.
\item $f_i=s_i\min(X,S)/S$; utility $\mathrm{E}[hf_i-\gamma f_i^2/2]$.
\item 100 lots: the size where the expected gain of a marginal lot equals its expected inventory cost.
\item Best responses to the others' total, damped, until they settle.
\item With many makers each one's share is small and a larger size buys share cheaply: the risk of exhausting the book never catches up before the
cap.
\item 1.1, 2.3, 6.2 and 50 (at the cap).
\item Lower at every $n$: 19.5 against 22.3 with ten makers, 12.3 against 14.3 with twenty.
\item Fills each of ten makers with 230 lots, 2.3 times the wanted size, all on the same side.
\item See \cref{def:hf:futures-market-making:implied}; 0 at a lag of one event, 0.9, 5.9, 23.7 and 49.7 per thousand at 5, 10, 20 and 30.
\item With ten makers each shows 2.3 times its wanted size, twenty makers go to the exchange's maximum; a 99.9th-percentile order leaves each maker
2.3 times its wanted inventory at once.
\item At the equilibrium size, with fast cancellation when large orders become likely, and a limit on the fill it can bear.
\item Top-order or time-weighted allocation, a minimum resting time, or charges on displayed but untraded size.
\item Consistency: every book's price comes from one curve, so the maker is never arbitraged between its own quotes.
\item When the relation between the two commodities changes (a refinery outage, a new delivery rule).
\item \omterm{def:hf:futures-market-making:implied}{Implied-price arbitrage}.
\item Model the others' sizes from the book's total depth and simulate the allocation with \texttt{firm.match}.
\item A share of every order, and with it the risk of every large one.
\end{enumerate}
\end{solution}

\begin{solution}{iq:hf:futures-market-making:1}
The tick is large relative to the contract's volatility over a trade's horizon, and many makers compete: the spread cannot be narrower than a tick,
and competition keeps it from being wider.
\iqlookfor{large tick, competition.}
\end{solution}

\begin{solution}{iq:hf:futures-market-making:2}
With $f_i=Xs_i/S$: $\mathrm{E}[f_i]=\mu s_i/S$, $\mathrm{E}[f_i^2]=m_2s_i^2/S^2$; maximise $h\mu p-\gamma m_2p^2/2$ over the share $p=s_i/S$:
$p^\ast=h\mu/(\gamma m_2)$, so $s_i=p^\ast S_{-i}/(1-p^\ast)$ if $p^\ast<1$. In a symmetric equilibrium $p=1/n$, so if $1/n<p^\ast$ every maker
wants more and sizes grow without bound: the book-exhaustion term is what stops it.
\iqlookfor{shares, and why the unbounded case arises.}
\end{solution}

\begin{solution}{iq:hf:futures-market-making:3}
You fill 2,000 times your share of the book; every maker is filled at once and the price likely moves through the level. Cancel all resting
orders on that side, hedge in the most liquid related contract, and reassess before quoting again.
\iqlookfor{correlated fills and immediate risk reduction.}
\end{solution}

\begin{solution}{iq:hf:futures-market-making:4}
Shares by floor of qty times size over total; drop shares below the minimum; leftovers by time. Edge cases: aggressor larger than the book, all
shares rounding to zero, orders fully filled before the leftover pass, ties in time.
\iqlookfor{rounding, minimums, exhaustion.}
\end{solution}

\begin{solution}{iq:hf:futures-market-making:5}
It is a breach if a full fill would exceed the limit: pre-trade risk checks must count displayed size as potential position; the firm needs a
limit on displayed size or a guarantee of cancellation.
\iqlookfor{potential position.}
\end{solution}

\begin{solution}{iq:hf:futures-market-making:6}
Compare depth with the distribution of aggressor sizes and with cancellations ahead of large prints; model makers' wanted sizes from their
behaviour in time-priority books of similar contracts.
\iqlookfor{depth versus executed size and cancellation behaviour.}
\end{solution}
